\begin{figure}[tb]
    \centering
    \tikzset{font=\fontfamily{phv}\selectfont}
    \resizebox{\columnwidth}{!}{%
    \begin{tikzpicture}
    \tikzstyle{instance} = [draw, minimum width=2.5cm, minimum height=1.5cm, ellipse, line width=0.8pt]
    \tikzstyle{type} = [draw, rounded corners, minimum width=2cm, minimum height=0.8cm, line width=0.8pt]
    \tikzstyle{arrow} = [->, thick,>=stealth']
    
    \tikzstyle{instanceLegend} = [draw, minimum width=1.8cm, minimum height=1cm, ellipse, line width=0.8pt]
    \tikzstyle{typeLegend} = [draw, rounded corners, minimum width=1.4cm, minimum height=0.6cm, line width=0.8pt]
    
    \matrix [draw,below left,row sep=0.3cm] at (11,-2) {%
      \node [instanceLegend] (p1){}; \node [right =0.3cm of p1] (p2){Instances}; \\
      \node [typeLegend] (p3){}; \node [right =0.55cm of p3] (p4){Types}; \\
    };
    
    % start node
    \node[instance, fill=tabl20lighgreen] (I1) {\textbf{Person 1}};
    \node[type, left = of I1, fill=tab20lightblue] (T1) {\textbf{Male}};
    \node[type, above left = of I1] (T2) {Grandfather};
    \node[type, above = of I1] (T3) {Father};
    \draw[arrow, line width=2pt] (I1) -- (T1);
    \draw[arrow] (I1) -- (T2);
    \draw[arrow] (I1) -- (T3);
    
    % married
    \node[instance, above right = 3cm of I1, , fill=tab20lightblue] (I2) {\textbf{Person 2}};
    \draw[arrow, line width=2pt] (I1) to  node[right=0.1cm, pos=0.4]  {\textbf{married}} (I2);
    \node[type, above = of I2] (T4) {Mother};
    \node[type, above left =1cm of I2] (T5) {Female};
    \draw[arrow] (I2) -- (T4);
    \draw[arrow] (I2) -- (T5);
    
    % married, hasSibling
    \node[instance, right = 3cm of I2, , fill=tab20lightblue] (I3) {\textbf{Person 3}};
    \draw[arrow, line width=2pt] (I2) to  node[below=0.1cm, pos=0.5]  {\textbf{hasSibling}} (I3);
    \node[type, above = of I3] (T6) {Female};
    \node[type, above left =1cm of I3, fill=tab20lightblue] (T7) {\textbf{Parent}};
    \draw[arrow] (I3) -- (T6);
    \draw[arrow, line width=2pt] (I3) -- (T7);
    
    % married, ...
    \node[instance, below = 1.75cm of I3] (I7) {\ldots};
    \draw[arrow] (I2) to  node[below=0.1cm, pos=0.5]  {\ldots} (I7);
    
    %hasParent
    \node[instance, right = 2cm of I1] (I4) {Person 4};
    \draw[arrow] (I1) to  node[below=0.1cm, pos=0.5]  {hasParent} (I4);
    
    % hasChild
    \node[instance, below = 2cm of I1, fill=tab20lightblue] (I5) {\textbf{Person 5}};
    \draw[arrow, line width=2pt] (I1) to  node[right=0.1cm, pos=0.6]  {\textbf{hasChild}} (I5);
    \node[type, left = of I5, fill=tab20lightblue] (T10) {\textbf{Child}};
    \node[type, above left = 1cm of I5] (T11) {Male};
    \draw[arrow, line width=2pt] (I5) -- (T10);
    \draw[arrow] (I5) -- (T11);
    
    % hasChild, ...
    \node[instance, right =2cm of I5] (I6) {\ldots};
    \draw[arrow] (I5) to  node[below=0.1cm, pos=0.5]  {\ldots} (I6);
    
    \end{tikzpicture}}
    \caption{Initialization of population: generating a concept via biased random walk (blue) originating from the chosen positive example (green).}
    \label{fig:init-example}
    \Description{Initialization of population: generating a concept via biased random walk originating from the chosen positive example.}
    \vspace{-1em}
\end{figure}
